\begin{theorem}[Nonderivability]
\label{thm:ch-nonderivability}
If \(\ZFC\) is consistent, then \(\ZFC\notderives\mathsf{CH}\) and
\(\ZFC\notderives\neg\mathsf{CH}\).
\end{theorem}

\begin{proof}
Suppose \(\ZFC\derives\mathsf{CH}\).  A witnessing derivation has its
open assumptions in \(\ZFC\subseteq\ZFC\cup\set{\neg\mathsf{CH}}\), so
monotonicity of the derivability relation gives
\(\ZFC+\neg\mathsf{CH}\derives\mathsf{CH}\); the assumption rule gives
\(\ZFC+\neg\mathsf{CH}\derives\neg\mathsf{CH}\); and negation elimination
gives \(\ZFC+\neg\mathsf{CH}\derives\bottom\).  The contrapositive of the
cited relative-consistency theorem for the negation of \(\mathsf{CH}\)
then refutes the consistency of \(\ZFC\), contrary to hypothesis.
Symmetrically, if \(\ZFC\derives\neg\mathsf{CH}\), then monotonicity, the
assumption rule, and negation elimination give
\(\ZFC+\mathsf{CH}\derives\bottom\), contradicting the consistency of
\(\ZFC+\mathsf{CH}\) supplied by \cref{thm:ch-relative-consistency}.
\end{proof}

